% from ega4-2-fr.tex
\oldpage[IV-2]{83}

\begin{proof}
En effet, il résulte de (4.8.8) qu'un tel corps est aussi un plus petit corps
de définition du graphe de $\overline u$, et l'existence de ce dernier résulte
de (4.8.9).
\end{proof}

\begin{corollary}[4.8.11]
\label{IV.2.4.8.11-fr}
Soient $k$ un corps, $X$ un $k$-préschéma, $K$ une extension de $k$,
$\overline Z$ un sous-préschéma fermé de $X_{(K)}$; alors il existe un plus
petit corps de définition de $\overline Z$.
\end{corollary}

\begin{proof}
En effet, si $\overline{\mathscr J}$ est l'Idéal quasi-cohérent de
$\mathcal O_{X_{(K)}}$ définissant $\overline Z$, il résulte de (4.8.8) que
cela équivaut à l'existence d'un plus petit corps de définition de
$\overline{\mathscr J}$, qui résulte de (4.8.9).
\end{proof}

\begin{corollary}[4.8.12]
\label{IV.2.4.8.12-fr}
Soient $k$ un corps, $X$, $Y$ deux $k$-préschémas, $Y$ étant supposé séparé,
$K$ une extension de $k$, $\overline f:X_{(K)}\to Y_{(K)}$ un $K$-morphisme.
Il existe alors un plus petit corps de définition de $\overline f$.
\end{corollary}

\begin{proof}
En effet, l'existence d'un tel corps, en vertu de (4.8.8), équivaut à celle
d'un plus petit corps de définition du graphe $\overline Z$ de $\overline f$;
mais comme $\overline Z$ est un sous-préschéma fermé de
$(X\times_kY)_{(K)}$ (I, 5.4.3), l'existence d'un plus petit corps de
définition de $\overline Z$ résulte de (4.8.11).
\end{proof}

\begin{proposition}[4.8.13]
\label{IV.2.4.8.13-fr}
Sous les hypothèses de (4.8.11) (resp. (4.8.9), (4.8.10), (4.8.12))
supposons en outre que $X$ soit de type fini sur $k$ (resp. que $X$ soit de
type fini sur $k$ et $\mathscr F$ cohérent, resp. que $X$ soit de type fini
sur $k$ et $\mathscr F$ et $\mathscr G$ cohérents, resp. que $X$ et $Y$
soient de type fini sur $k$). Alors le plus petit corps de définition de
$\overline Z$ (resp. de $\overline{\mathscr H}$, de $\overline u$, de $\overline f$) est
une extension de type fini de $k$.
\end{proposition}

\begin{proof}
On peut se borner à traiter le cas de (4.8.9), et supposer $X$ affine puisque
$X$ est réunion finie d'ouverts affines de type fini sur $k$. Avec les
notations de (4.8.6), tout revient à montrer que si $A$ est une algèbre de
type fini sur $k$ et $M$ un $A$-module de type fini, il y a une sous-extension
$K'$ de $K$ vérifiant les conditions a), b) ou c) et qui soit de type fini sur
$k$ (toute sous-extension d'une extension de type fini étant de type fini).
Or, soit $(m_i)_{1\leq i\leq n}$ un système de générateurs du $A$-module $M$;
comme $A$ est noethérien, il en est de même de $A_{(K)}$, donc $\overline N$
est un $A_{(K)}$-module de type fini, admettant un système fini de générateurs
$n_j$ ($1\leq j\leq r$) tels que
\[
n_j=\sum_i(m_i\otimes1)b_{ij},
\quad\text{où }b_{ij}\in A_{(K)};
\]
chaque $b_{ij}$ s'écrit d'ailleurs
\[
b_{ij}=\sum_h a_{ijh}\otimes c_{ijh},
\quad\text{où }a_{ijh}\in A\text{ et }c_{ijh}\in K;
\]
il est clair que l'extension $K'$ de $k$ engendrée par les $c_{ijh}$ répond à
la question.
\end{proof}

\begin{proposition}[4.8.14]
\label{IV.2.4.8.14-fr}
Soient $k$ un corps, $X$ un $k$-préschéma de type fini, $K$ une extension de
$k$ contenant une extension parfaite de $k$. Alors le plus petit corps de
définition du sous-préschéma fermé $(X_{(K)})_{\mathrm{red}}$ de $X_{(K)}$
est une extension finie radicielle de $k$.
\end{proposition}

\begin{proof}
Si $p$ est l'exposant caractéristique de $k$, l'hypothèse entraîne que
$k'=k^{p^{-\infty}}$ est contenu dans $K$; on sait que
$(X_{(k')})_{\mathrm{red}}$ est géométriquement réduit sur $k'$ (4.6.1), donc
$(X_{(k')})_{\mathrm{red}}\otimes_{k'}K$ est réduit, et par suite égal à
$(X_{(K)})_{\mathrm{red}}$ (I, 5.1.8); en d'autres termes, $k'$ est un corps
de définition de $(X_{(K)})_{\mathrm{red}}$. Comme $k'$ est une extension
algébrique de $k$, la conclusion résulte de (4.8.13).
\end{proof}

\begin{remarks}[4.8.15]
\label{IV.2.4.8.15-fr}
\begin{enumerate}
\item[(i)] La conclusion de (4.8.14) ne subsiste plus nécessairement si l'on
ne suppose pas que $K$ contient une extension parfaite de $k$. Par exemple,
soient $k_0$

